% =============================================================================
% INSTRUCTION SWITCH DATASET (ISD) - CITA Paper in ALKALI Format
% =============================================================================

\vspace{-2mm}
\section{Instruction Switch Dataset (ISD)}
\label{sec:isd_dataset}

To rigorously evaluate instruction-conditioned alignment, we construct the \textbf{Instruction Switch Dataset (ISD)}---a diagnostic benchmark where the \textbf{user prompt remains constant} while only the \textbf{alignment instruction varies}. This design isolates the effect of instructions from prompt-specific confounds.

\vspace{-2mm}
\subsection{Motivation}
\label{sec:isd_motivation}

Existing benchmarks conflate instruction-following with prompt variation. When both change simultaneously, it is impossible to determine whether behavioral differences stem from instruction comprehension or prompt-specific responses.

\textbf{ISD Design Principle.} By fixing the user prompt $X$ and varying only the instruction $I$, we can measure:
\begin{equation}
\Delta_I = \mathbb{E}_{I_1, I_2}\left[ d\left( R(I_1, X), R(I_2, X) \right) \right]
\label{eq:isd_metric}
\end{equation}

where $R(I, X)$ is the response to prompt $X$ under instruction $I$, and $d(\cdot, \cdot)$ measures behavioral difference.

\vspace{-2mm}
\subsection{Dataset Construction}
\label{sec:isd_construction}

\paragraph{Prompt Bank.} We curate 500 prompts from diverse sources:
\begin{itemize}[leftmargin=*,itemsep=1pt]
    \item News articles and current events
    \item Educational Q\&A collections
    \item Instruction-tuning datasets (FLAN~\cite{wei2022finetuned}, Alpaca~\cite{taori2023alpaca})
    \item User intent pools from dialogue systems
\end{itemize}

\paragraph{Selection Criteria.} All prompts are:
\begin{itemize}[leftmargin=*,itemsep=1pt]
    \item \textbf{Semantically open-ended}: Allow multiple valid response styles
    \item \textbf{Culturally grounded}: Avoid region-specific assumptions
    \item \textbf{Neutrally phrased}: No embedded biases toward specific instruction types
\end{itemize}

\vspace{-2mm}
\subsection{10 Instruction Types}
\label{sec:instruction_types}

\begin{table}[h!]
\centering
\small
\begin{tabular}{@{}ll@{}}
\toprule
\textbf{Instruction Type} & \textbf{Behavioral Objective} \\
\midrule
\texttt{neutral} & Balanced perspective, present pros/cons \\
\texttt{conservative} & Favor traditional methods, risk-averse \\
\texttt{liberal} & Encourage innovation, embrace change \\
\texttt{regulatory} & Emphasize compliance, legal considerations \\
\texttt{empathetic} & Compassionate, emotionally supportive tone \\
\texttt{safety\_first} & Prioritize risk awareness, caution \\
\texttt{educational} & Teaching mode, explain concepts clearly \\
\texttt{concise} & Brief, direct responses \\
\texttt{professional} & Formal tone, business-appropriate \\
\texttt{creative} & Novel approaches, unconventional ideas \\
\bottomrule
\end{tabular}
\caption{The 10 instruction types in ISD, each targeting a distinct behavioral dimension.}
\label{tab:instruction_types}
\vspace{-4mm}
\end{table}

\vspace{-2mm}
\subsection{Dataset Scale}
\label{sec:isd_scale}

\begin{equation}
|\text{ISD}| = 300 \text{ prompts} \times 10 \text{ instructions} = \mathbf{3,000} \text{ test cases}
\label{eq:isd_scale}
\end{equation}

\vspace{-2mm}
\subsection{Example: Instruction Switching}
\label{sec:isd_example}

\textbf{Prompt}: ``\textit{Should AI-generated art be allowed in competitions?}''

\begin{table}[h!]
\centering
\footnotesize
\begin{tabular}{@{}lp{5.5cm}@{}}
\toprule
\textbf{Instruction} & \textbf{Response Snippet} \\
\midrule
Neutral & AI art offers exciting possibilities but raises authenticity concerns that deserve careful consideration. \\
Conservative & Competitions should prioritize human craftsmanship and traditional artistic skills. \\
Liberal & AI art promotes innovation, democratizes creativity, and expands artistic boundaries. \\
Regulatory & Ensure copyright compliance, require disclosure, and label AI-generated entries. \\
Empathetic & We must respect artists' concerns about their livelihoods while embracing new tools. \\
\bottomrule
\end{tabular}
\caption{Same prompt, different instructions produce behaviorally distinct responses.}
\label{tab:isd_example}
\vspace{-4mm}
\end{table}

\vspace{-2mm}
\subsection{Prompt Category Distribution}
\label{sec:prompt_categories}

\begin{table}[h!]
\centering
\small
\begin{tabular}{@{}lc|lc@{}}
\toprule
\textbf{Category} & \textbf{\#} & \textbf{Category} & \textbf{\#} \\
\midrule
Technology & 25 & Ethics & 25 \\
Healthcare & 25 & Culture & 25 \\
Environment & 25 & Science & 25 \\
Education & 25 & Business & 25 \\
Economics & 25 & Personal & 25 \\
Social & 25 & Governance & 25 \\
\bottomrule
\end{tabular}
\caption{ISD prompt distribution across 12 categories (300 prompts total).}
\label{tab:prompt_categories}
\vspace{-4mm}
\end{table}

\vspace{-2mm}
\subsection{Instruction-Following vs. Instruction-Alignment}
\label{sec:following_vs_alignment}

A critical distinction underlies the ISD evaluation:

\begin{table}[h!]
\centering
\small
\begin{tabular}{@{}lcc@{}}
\toprule
\textbf{Property} & \textbf{Following} & \textbf{Alignment} \\
\midrule
Depth & Surface-level & Internalized \\
Example & ``Be concise'' $\rightarrow$ short text & Concise \textit{style} adopted \\
Robustness & Easily disrupted & Consistent across contexts \\
Goal & Format compliance & Policy embodiment \\
\bottomrule
\end{tabular}
\caption{Instruction-following (shallow) vs. instruction-alignment (deep).}
\label{tab:following_vs_alignment}
\vspace{-4mm}
\end{table}

\textbf{ISD tests whether CITA achieves true instruction-alignment}---consistent behavioral shifts reflecting internalized policy goals---rather than superficial instruction-following that can be easily circumvented.

\vspace{-2mm}
\subsection{Dataset Availability}
\label{sec:dataset_availability}

The Instruction Switch Dataset is publicly available:

\begin{center}
\url{https://huggingface.co/datasets/kapilw25/ISD-Instruction-Switch-Dataset}
\end{center}
